Bagian ini membahas konsep-konsep teoritis yang menjadi landasan pengembangan
program OctoVox. Pemahaman terhadap konsep-konsep berikut diperlukan untuk
memahami cara kerja algoritma \textit{voxelization} berbasis \textit{Octree} yang
diimplementasikan.

\section{Format File \textit{.obj}}

Format \textit{Wavefront OBJ} (.obj) adalah format file teks yang digunakan
secara luas untuk merepresentasikan geometri objek tiga dimensi. Format ini
bersifat sederhana dan dapat dibaca langsung oleh manusia karena tersimpan dalam
bentuk teks biasa. Setiap baris dalam file \texttt{.obj} mendefinisikan satu
elemen geometri dengan awalan tertentu.

Dalam konteks program OctoVox, hanya dua jenis elemen yang diproses:

\begin{itemize}
    \item \textbf{Vertex} (\texttt{v}): mendefinisikan sebuah titik dalam ruang
    tiga dimensi dengan koordinat $x$, $y$, dan $z$. Setiap \textit{vertex}
    secara implisit memiliki nomor urut yang dimulai dari 1 sesuai urutan
    kemunculannya dalam file.

    \item \textbf{Face} (\texttt{f}): mendefinisikan sebuah poligon dengan
    mereferensikan indeks \textit{vertex} yang membentuknya. Pada program ini,
    setiap \textit{face} diasumsikan berbentuk segitiga sehingga hanya terdiri
    dari tiga indeks \textit{vertex}.
\end{itemize}

Kumpulan \textit{face} yang saling terhubung membentuk apa yang disebut sebagai
\textit{polygon mesh}, yaitu representasi permukaan suatu objek tiga dimensi
sebagai jaring-jaring segitiga yang saling berbagi \textit{vertex} dan sisi.

\section{Representasi Ruang Tiga Dimensi dan \textit{Bounding Box}}

\subsection{Sistem Koordinat Kartesius Tiga Dimensi}

Ruang tiga dimensi dalam komputasi grafis direpresentasikan menggunakan sistem
koordinat Kartesius dengan tiga sumbu orthogonal, yaitu sumbu $x$, $y$, dan $z$.
Setiap titik dalam ruang tiga dimensi didefinisikan oleh tiga buah bilangan riil
yang masing-masing menyatakan posisi titik tersebut terhadap ketiga sumbu.

\subsection{\textit{Axis-Aligned Bounding Box}}

\textit{Axis-Aligned Bounding Box} (AABB) adalah kotak pembatas terkecil yang
mengelilingi suatu objek atau sekumpulan titik dalam ruang tiga dimensi, di mana
sisi-sisinya sejajar dengan sumbu-sumbu koordinat. AABB didefinisikan oleh dua
buah titik, yaitu titik minimum $P_{\min} = (x_{\min}, y_{\min}, z_{\min})$ dan
titik maksimum $P_{\max} = (x_{\max}, y_{\max}, z_{\max})$, yang masing-masing
merepresentasikan pojok terdekat dan pojok terjauh dari pusat koordinat.

Nilai $P_{\min}$ dan $P_{\max}$ diperoleh dengan mencari nilai minimum dan
maksimum dari koordinat seluruh \textit{vertex} objek pada setiap sumbu:

\begin{align}
    x_{\min} &= \min_{i} x_i, \quad x_{\max} = \max_{i} x_i \\
    y_{\min} &= \min_{i} y_i, \quad y_{\max} = \max_{i} y_i \\
    z_{\min} &= \min_{i} z_i, \quad z_{\max} = \max_{i} z_i
\end{align}

Dalam konteks \textit{voxelization} berbasis \textit{Octree}, AABB objek
dinormalisasi menjadi sebuah kubus sempurna dengan mengambil panjang sisi
terpanjang sebagai ukuran sisi kubus, sehingga seluruh objek terkandung di dalam
kubus tersebut.

\section{Voksel dan \textit{Voxelization}}

\subsection{Voksel}

Voksel (\textit{voxel}), yang merupakan kependekan dari \textit{volumetric pixel},
adalah unit terkecil dalam representasi tiga dimensi berbasis volume. Voksel
merepresentasikan sebuah nilai pada suatu posisi reguler dalam ruang tiga dimensi,
analog dengan piksel pada gambar dua dimensi. Dalam konteks program ini, voksel
direpresentasikan sebagai kubus berukuran seragam yang menempati suatu posisi
dalam ruang tiga dimensi.

\subsection{\textit{Voxelization}}

\textit{Voxelization} adalah proses konversi representasi geometri kontinu, seperti
\textit{polygon mesh}, menjadi representasi diskrit berbasis voksel. Terdapat dua
jenis \textit{voxelization} yang umum dikenal:

\begin{itemize}
    \item \textbf{\textit{Surface voxelization}}: hanya menghasilkan voksel pada
    bagian yang bersinggungan dengan permukaan objek. Jenis inilah yang
    diimplementasikan dalam OctoVox.

    \item \textbf{\textit{Solid voxelization}}: menghasilkan voksel pada seluruh
    ruang interior objek, tidak hanya pada permukaannya.
\end{itemize}

Resolusi hasil \textit{voxelization} ditentukan oleh ukuran voksel yang digunakan.
Semakin kecil ukuran voksel, semakin tinggi resolusi dan semakin detail objek yang
direpresentasikan, namun semakin besar pula kebutuhan memori dan komputasi.

\section{\textit{Octree}}

\textit{Octree} adalah struktur data pohon hierarkis yang digunakan untuk
mempartisi ruang tiga dimensi secara rekursif. Setiap \textit{node} dalam
\textit{Octree} merepresentasikan sebuah kubus dalam ruang tiga dimensi. Apabila
suatu \textit{node} perlu diperhalus lebih lanjut, \textit{node} tersebut dibagi
menjadi tepat delapan \textit{node} anak yang merepresentasikan delapan sub-kubus
berukuran sama hasil pembagian kubus induknya. Pembagian ini dilakukan dengan
membelah kubus tepat di tengah pada ketiga sumbu $x$, $y$, dan $z$ secara
bersamaan.

\subsection{Terminologi}

Beberapa istilah penting dalam \textit{Octree} adalah sebagai berikut:

\begin{itemize}
    \item \textbf{\textit{Root node}}: \textit{node} paling atas dalam hierarki
    \textit{Octree} yang merepresentasikan seluruh ruang yang diperhitungkan,
    dalam konteks ini adalah AABB objek.

    \item \textbf{\textit{Internal node}}: \textit{node} yang memiliki delapan
    \textit{node} anak dan merepresentasikan suatu sub-ruang yang masih perlu
    diperhalus.

    \item \textbf{\textit{Leaf node}}: \textit{node} yang tidak memiliki anak,
    baik karena telah mencapai kedalaman maksimum maupun karena sub-ruang yang
    diwakilinya dinyatakan kosong.

    \item \textbf{Kedalaman (\textit{depth})}: jarak suatu \textit{node} dari
    \textit{root node}. \textit{Root node} berada pada kedalaman 0.

    \item \textbf{Oktan}: salah satu dari delapan sub-kubus hasil pembagian suatu
    \textit{node}.
\end{itemize}

\subsection{\textit{Pruning}}

\textit{Pruning} adalah mekanisme pemangkasan cabang \textit{Octree} yang tidak
perlu ditelusuri lebih lanjut. Dalam konteks \textit{voxelization}, sebuah
\textit{node} dapat dipangkas apabila sub-ruang yang diwakilinya terbukti tidak
bersinggungan dengan permukaan objek. Mekanisme ini secara signifikan mengurangi
jumlah \textit{node} yang perlu diproses, sehingga meningkatkan efisiensi
komputasi secara keseluruhan.

Jumlah \textit{node} pada kedalaman $d$ tanpa \textit{pruning} adalah $8^d$,
sehingga total \textit{node} pada kedalaman maksimum $N$ tanpa \textit{pruning}
adalah:
\begin{equation}
    \sum_{d=0}^{N} 8^d = \frac{8^{N+1} - 1}{7}
\end{equation}

Dengan adanya \textit{pruning}, hanya \textit{node} yang bersinggungan dengan
permukaan objek yang ditelusuri, sehingga jumlah \textit{node} yang diproses
jauh lebih kecil dari nilai tersebut.

